% ============================================================================
%  Capitolo 7 -- Verm1, parte 1: una testa che attraversa il mondo (Lezione 07)
% ============================================================================
\chapter{Verm1, parte 1: una testa che attraversa il mondo}
\label{ch:verm1p1}

\epigraphline{Un viaggio di mille miglia comincia con un solo passo.}
  {Laozi, \emph{Tao Te Ching}}

\begin{chapgoals}
Alla fine di questo capitolo avrete
\begin{itemize}
\item una cartuccia vostra, \fname{verm1.prg32}, costruita da un solo file
  assembly scritto da voi;
\item un titolo, una griglia di 40 per 24 celle, e la testa del verme disegnata in una di esse;
\item la testa che si muove da sola, una cella alla volta, a una velocità che controllate in virgola fissa;
\item un mondo senza bordi: oltre l'ultima colonna la testa rientra dalla prima;
\item l'abitudine al ciclo di lavoro: modificare, compilare, eseguire,
  guardare, e fermare il programma nel debugger.
\end{itemize}
\end{chapgoals}
% Generated by tools/textbook/export_listings.js from tools/decks/mastery.js and tools/concepts/registry.json: do not edit.
\begin{chapstatus}{Stato di apprendimento di questo capitolo}
\masterychip{now}{DA COMPRENDERE OGGI}\enspace il progetto Verm1; contratto init/update/draw; un colore è un numero; gli array in memoria; .text, .rodata, .data; indirizzo, valore, load, store; perché si salva ra; shift come moltiplicazione; la virgola fissa; le pseudo-istruzioni; l'ordine little-endian.\par\smallskip
\masterychip{use}{DA UTILIZZARE OGGI}\enspace gli strumenti di PRG32, il debugger di PRG32-QT (si usa, non si studia, in questo corso).\par\smallskip
\masterychip{later}{DA COMPRENDERE PIÙ AVANTI}\enspace formati delle istruzioni (Capitolo~\ref{ch:instructions}).
\end{chapstatus}
\begin{chapstatus}{Percorso del capitolo}
Alla fine del capitolo saprete:
\begin{enumerate}[nosep,leftmargin=2em]
\item costruire una cartuccia da un file vuoto: una griglia, un titolo e una testa che si muove
\item mettere costanti e variabili nella sezione giusta, .rodata o .data
\item controllare una velocità con un accumulatore in virgola fissa
\end{enumerate}
Si appoggia su: contratto init/update/draw (Capitolo~\ref{ch:firstcontact}); un colore è un numero (Capitolo~\ref{ch:information}); gli array in memoria (Capitolo~\ref{ch:information}); indirizzo, valore, load, store (Capitolo~\ref{ch:firstcontact}); perché si salva ra (Capitolo~\ref{ch:firstcontact}). Prossimo passo: Capitolo~\ref{ch:boole1}.
\par\smallskip\textbf{Guardando indietro, guardando avanti.} Nella Lezione 02 il passo 2 salvava ra sullo stack prima delle sue chiamate, e lw t1, 0(t0) leggeva una parola all'indirizzo in t0. Oggi verm1\_draw salva ra e verm1\_update no, e ogni variabile di Verm1 si legge con la e lw. La Lezione 18 spiega lo stack e la convenzione di chiamata; la Lezione 19 nomina i modi di indirizzamento. \emph{Richiamo:} Perché verm1\_draw salva ra mentre verm1\_update no?
\end{chapstatus}


Questo è il primo tutorial del progetto. Usa soltanto ciò che i Capitoli da
\ref{ch:firstcontact} a~\ref{ch:reals} hanno insegnato, ed è pensato per essere
seguito alla tastiera: il programma lo scrivete voi. Tutto ciò che è stampato
nei listati è letto dal file di riferimento \fname{verm1\_part1.S}, che viene
compilato ed eseguito automaticamente; il file intero è
nell'Appendice~\ref{app:verm1}.

% ----------------------------------------------------------------------------
\section{Il progetto}

\subsection{Verm1: un omaggio a Snake}

\vermone, che si legge \emph{vermone}, è un verme che si muove su una griglia,
mangia, cresce, e non deve mordersi. Il gioco ha una lunga storia: apparve come
\emph{Blockade}, una macchina da sala giochi della Gremlin, nel 1976, e come
\emph{Snake} sul Nokia~6110 nel 1997, scritto da Taneli Armanto, divenne uno
dei giochi più giocati di sempre. \vermone è un omaggio, non una copia: il
nome, le regole di dettaglio e l'aspetto sono suoi. È abbastanza piccolo da
essere scritto da zero in assembly e abbastanza ricco da richiedere ogni
argomento del corso, e ogni sua riga la scrivete voi: non c'è un motore, e non
c'è una libreria oltre le chiamate di \prg. L'implementazione di riferimento
che i capitoli tutorial seguono, e con la quale potete sempre confrontare un
programma vostro, è un contributo di Giuseppe Coviello
(Appendice~\ref{app:verm1}). La Figura~\ref{fig:v1:goal} mostra la
destinazione.

\begin{figure}[htb]
\centering
\includegraphics[width=0.55\textwidth]{verm1_p4_play.png}
\caption{\vermone alla fine del libro (parte 4).}\label{fig:v1:goal}
\end{figure}

\subsection{Il piano}

\begin{table}[htb]
\centering\small
\caption{Quattro parti e una pubblicazione.}\label{tab:v1:plan}
\begin{tabularx}{\textwidth}{@{}lLl@{}}
\toprule
\textbf{Capitolo} & \textbf{Aggiungete} & \textbf{Usa i Capitoli}\\
\midrule
\ref{ch:verm1p1}\quad Parte 1 & una griglia, una testa, il movimento automatico, il rientro dal bordo opposto & da \ref{ch:history} a \ref{ch:reals}: informazione e numeri\\
\ref{ch:verm1p2}\quad Parte 2 & la guida con il joystick, il cibo, un punteggio & da \ref{ch:boole1} a \ref{ch:combinational}: logica e circuiti\\
\ref{ch:verm1p3}\quad Parte 3 & un corpo che cresce, le collisioni, le schermate del titolo e di fine partita & da \ref{ch:sequential} a \ref{ch:storage}: stato e memoria\\
\ref{ch:verm1p4}\quad Parte 4 & il disegno veloce, il punteggio in cifre, i livelli, il suono & da \ref{ch:processor} a \ref{ch:io}: la macchina\\
\ref{ch:store}\quad Lo Store & un bundle con manifesto, icona e crediti & tutti\\
\bottomrule
\end{tabularx}
\end{table}

Ogni parte termina con un programma che funziona (Tabella~\ref{tab:v1:plan}).

\begin{prgwarn}
Non cominciate mai una parte con un programma che non funziona. Se il vostro è
rotto, ripartite dal file di riferimento della parte precedente.
\end{prgwarn}

La parte~1 prende dai capitoli letti finora ciò che segue.

\begin{center}\small
\begin{tabularx}{\textwidth}{@{}lLL@{}}
\toprule
\textbf{Da} & \textbf{Idea} & \textbf{Dove compare}\\
\midrule
Capitolo \ref{ch:firstcontact} & i registri, \ins{la}/\ins{lw}/\ins{sw}, \ins{call} e lo stack frame & dappertutto\\
Capitolo \ref{ch:history} & i giochi retro come esemplari & la scelta del gioco\\
Capitolo \ref{ch:information} & i colori come numeri RGB565; le stringhe terminate da zero & \code{GREEN}, il titolo\\
Capitolo \ref{ch:integers} & gli scorrimenti come moltiplicazioni; i confronti; il rientro dal bordo & da cella a pixel; il bordo del mondo\\
Capitolo \ref{ch:reals} & la virgola fissa: una frazione come intero scalato & la velocità del verme\\
\bottomrule
\end{tabularx}
\end{center}

% ----------------------------------------------------------------------------
\section{Passo 0: il ciclo di lavoro}

Create una cartella \fname{\textasciitilde/verm1} e teneteci un solo file,
\fname{verm1.S}, che crescerà di parte in parte. I due comandi del
Listato~\ref{lst:v1:build} compilano la cartuccia e la inviano a \prgqt; sono
gli stessi in ogni tutorial. \prgqt ascolta sulla porta 8080 e mostra il suo
indirizzo nella schermata Setup: usate quell'indirizzo se \code{localhost} non
risponde. La stessa cartuccia gira su QEMU (\code{qemu upload}, poi
\code{qemu run}) e sulla scheda (\code{-{}-url http://192.168.4.1}).

\shfile[caption={Compilare ed eseguire: modificare, compilare, eseguire, guardare.},label=lst:v1:build]{code/l07/c03.sh}

% ----------------------------------------------------------------------------
\section{Passo 1: tre funzioni che ritornano}

\asmfile[caption={Passo 1 \textbf{(nuovo)}: lo scheletro.},label=lst:v1:skel]{code/l07/c04.S}

È lo scheletro del Capitolo~\ref{ch:firstcontact} con un altro prefisso:
\code{verm1} è il nome passato a \code{-{}-entry-prefix}. \code{init} viene
eseguita una volta; \code{update} e \code{draw} una volta per frame.

\begin{prgcheck}[Eseguite]
Scrivetelo, compilatelo, eseguitelo: uno schermo nero è la risposta giusta. Se
qui la compilazione fallisce il problema è l'installazione, non il programma:
risolvetelo prima di proseguire. D'ora in poi: aggiungete una cosa, eseguite,
guardate. Mai due cose insieme.
\end{prgcheck}

% ----------------------------------------------------------------------------
\section{Passo 2: una griglia e una testa}

\subsection{Costanti: i colori e la griglia}

Sostituite le prime righe dello scheletro con quelle del
Listato~\ref{lst:v1:const}.

\asmfile[caption={Passo 2 \textbf{(modificato)}: l'inizio del file.},label=lst:v1:const]{code/l07/c05.S}

I colori sono numeri RGB565: \hex{07e0} è verde puro. Lo schermo è largo 320
pixel e una cella è larga 8 pixel, il che fa 40 colonne. Gli 8 pixel in alto
contengono il titolo, e ne restano 192: 24 righe. La direttiva \code{.equ} dà
un nome a un numero e non occupa memoria. La costante \code{SPEED\_Q8} si usa
nel passo~3.

\subsection{Il mondo di Verm1}

\begin{prgconcept}
Il verme non vive nei pixel: vive nelle celle. Una posizione è una colonna, da
0 a 39, e una riga, da 0 a 23.
\end{prgconcept}

Per disegnare una cella le sue coordinate vengono trasformate in pixel
(Figura~\ref{fig:v1:grid}): $x = \mathit{colonna}\times 8$, uno scorrimento a
sinistra di 3, e $y = \mathit{riga}\times 8+8$, perché prima viene la riga del
titolo. La cella $(0,0)$ è al pixel $(0,8)$; la cella $(39,23)$ è a
$(312,192)$. Pensare in celle rende semplice ogni regola successiva: muoversi è
sommare~1.

\begin{figure}[htb]
\centering
\begin{tikzpicture}[x=2.6mm,y=-2.6mm]
  \fill[prgsteel] (0,0) rectangle (40,1);
  \fill[prgink] (0,1) rectangle (40,25);
  \draw[prgsteel,very thin,step=1] (0,1) grid (40,25);
  \fill[green!80!black] (5,13) rectangle (6,14);
  \node[text=white,font=\sffamily\tiny,anchor=west] at (0.3,0.5) {riga del titolo: 8 pixel};
  \node[lbl] at (20,-1) {colonne $0\ldots 39$\quad ($x=\mathit{colonna}\times 8$)};
  \node[lbl,rotate=90] at (-3.2,13) {righe $0\ldots 23$\quad ($y=\mathit{riga}\times 8+8$)};
  \node[lbl] at (5.5,26) {colonna 5};
  \node[lbl,anchor=east] at (-0.3,13.5) {riga 12};
\end{tikzpicture}
\caption{Il mondo di \vermone: 40 per 24 celle di 8 per 8 pixel. La testa
parte dalla colonna 5 della riga 12.}
\label{fig:v1:grid}
\end{figure}

\subsection{Variabili e una stringa}

Mettete il Listato~\ref{lst:v1:data} alla fine del file. La sezione
\code{.rodata} contiene le costanti: il titolo termina con un byte a zero. La
sezione \code{.data} contiene le variabili, una parola di 32 bit ciascuna. I
loro valori iniziali sono solo un punto di partenza: \code{init} li imposterà
di nuovo.

\asmfile[caption={Passo 2 \textbf{(nuovo)}: i dati, alla fine del file.},label=lst:v1:data]{code/l07/c06.S}

\subsection{\texttt{init}: mettere la testa al suo posto}

\asmfile[caption={Passo 2 \textbf{(modificato)}: \code{verm1\_init}.},label=lst:v1:init]{code/l07/c07.S}

\ins{la} prende un \emph{indirizzo}, \ins{sw} vi memorizza un \emph{valore}. La
funzione imposta la partenza in modo esplicito: un programma non dovrebbe mai
fidarsi di ciò che la memoria conteneva. \reg{zero} è il registro che si legge
sempre 0. La funzione è una foglia: non fa chiamate, quindi non le serve uno
stack frame.

\subsection{\texttt{draw}: il titolo e la testa}

\asmfile[caption={Passo 2 \textbf{(modificato)}: \code{verm1\_draw}.},label=lst:v1:draw]{code/l07/c08.S}

Cancellare lo schermo, scrivere il titolo, disegnare un rettangolo: tre
chiamate, con gli argomenti che viaggiano da \reg{a0} ad \reg{a4}, in ordine.
Le firme sono
\code{prg32\_gfx\_text8(x, y, stringa, primo piano, sfondo)} e
\code{prg32\_gfx\_rect(x, y, larghezza, altezza, colore)}. \ins{slli} di 3
moltiplica per 8: da cella a pixel. Poiché \code{draw} fa chiamate, salva
\reg{ra} in uno stack frame.

\begin{prgcheck}[Eseguite: una testa sulla griglia]
Risultato atteso ora: la parola \texttt{VERM1} in alto a sinistra e un quadrato
verde fermo in colonna 5, riga 12 (la Figura~\ref{fig:v1:shot} lo mostra più
tardi, quando ha cominciato a muoversi).
\end{prgcheck}

\begin{figure}[htb]
\centering
\includegraphics[width=0.55\textwidth]{verm1_p1.png}
\caption{Il titolo e la testa. Questa immagine è stata presa dai test
automatici 64 frame dopo l'avvio della parte~1 finita: la testa è in
colonna 15.}
\label{fig:v1:shot}
\end{figure}

% ----------------------------------------------------------------------------
\section{Passo 3: il movimento}

\subsection{Quanto deve andare veloce un verme?}

\code{update} viene eseguita circa trenta volte al secondo. Una cella per frame
attraverserebbe lo schermo in poco più di un secondo: ingiocabile. Ciò che
serve è una frazione di cella per frame, e il processore del kit ha solo
interi. La risposta è la virgola fissa del Capitolo~\ref{ch:reals}, nel formato
Q24.8: il programma tiene una frazione in unità di $1/256$ di cella, e fa un
passo ogni volta che si è accumulata una cella intera. Con
$\code{SPEED\_Q8}=40$ il verme avanza di $40/256\approx 0,156$ celle per
frame: circa un passo ogni 6 frame.

È la palla del Capitolo~\ref{ch:reals} trasformata in un timer. Il resto non
viene buttato via, quindi la velocità media è esatta: 256 frame danno
esattamente 40 passi.

\subsection{\texttt{update}: accumulare, fare un passo, rientrare}

\asmfile[caption={Passo 3 \textbf{(modificato)}: \code{verm1\_update}.},label=lst:v1:update]{code/l07/c10.S}

La velocità viene sommata alla frazione; sotto 256, la funzione aspetta. A 256
sottrae 256, conserva il resto, e sposta la testa di una cella a destra. La
colonna 40 non esiste: oltre il bordo destro la testa torna alla colonna~0. La
funzione ha due uscite, e ciascuna memorizza ciò che ha cambiato. \ins{blt}
confronta numeri con segno, il che qui è corretto perché ogni valore è piccolo
e positivo. Le etichette locali \code{1:} e \code{2:} si richiamano come
\code{1f} e \code{2f}: «il prossimo 1, o 2, \emph{in avanti}» (\emph{forward}).

\subsection{Tracciatela a mano}

\begin{table}[htb]
\centering\small
\caption{Una traccia di \code{verm1\_update}.}\label{tab:v1:trace}
\begin{tabular}{@{}cccccc@{}}
\toprule
\textbf{Frame} & \textbf{\code{step\_q8} prima} & \textbf{+ 40} & \textbf{Cella intera?} & \textbf{\code{step\_q8} dopo} & \textbf{\code{head\_x}}\\
\midrule
1 & 0 & 40 & no & 40 & 5\\
2 & 40 & 80 & no & 80 & 5\\
6 & 200 & 240 & no & 240 & 5\\
7 & 240 & 280 & sì: $280-256$ & 24 & 6\\
13 & 224 & 264 & sì: $264-256$ & 8 & 7\\
64 & -- & -- & -- & 0 & 15\\
\bottomrule
\end{tabular}
\end{table}

I passi arrivano ogni 6 o 7 frame (Tabella~\ref{tab:v1:trace}): il resto porta
con sé la differenza, come fa un anno bisestile. Dopo 64 frame la somma è
$64\times 40 = 2560 = 10\times 256$: dieci passi, e la testa è in colonna 15,
esattamente come nella Figura~\ref{fig:v1:shot}. Completate i frame da 8 a 12
per esercizio.

\begin{prgcheck}[Eseguite: la testa attraversa il mondo]
\begin{enumerate}[nosep]
\item Compilate ed eseguite: la testa si muove verso destra, circa cinque celle al secondo.
\item Aspettate il bordo destro: la testa esce dalla colonna 39 e rientra dalla
  colonna 0.
\item Cambiate \code{SPEED\_Q8} in 8, poi in 128, poi in 256: prevedete la
  velocità prima di eseguire.
\item Cambiate la riga di partenza in \code{init}, e il colore della testa.
\end{enumerate}
\textbf{Risultato atteso.} La testa attraversa le quaranta colonne in 256
frame, circa otto secondi e mezzo a trenta frame al secondo, per sempre. Con
$\code{SPEED\_Q8}=256$ si muove di una cella per frame; valori più grandi danno
comunque un passo per frame, perché \code{update} fa al più un passo: un limite
che vale la pena scoprire.
\end{prgcheck}

% ----------------------------------------------------------------------------
\section{Passo 4: guardare dentro con PRG32-QT}

Il debugger fa parte del metodo, non è uno strumento di emergenza.

\begin{prgtry}[Fermate il programma e leggete la macchina]
\begin{enumerate}[nosep]
\item In \prgqt aprite \emph{Settings} e abilitate \emph{Run cartridges in
  RISC-V debug mode}.
\item Eseguite la vostra cartuccia: il gioco è a sinistra, il debugger a destra.
\item Premete \key{Update}: il listato salta a \code{verm1\_update}. Fate clic
  sul margine per mettere un breakpoint sulla sua prima istruzione.
\item Il programma si ferma lì al frame successivo. Premete \key{Step} e
  guardate i registri cambiare, un'istruzione alla volta.
\item Dopo le due istruzioni di \ins{la}, \reg{t0} contiene un indirizzo:
  scrivetelo nel campo \emph{Memory} per vedere \code{step\_q8} in memoria.
\item Premete \key{Resume}, e guardate i byte di \code{step\_q8} crescere di 40
  (\hex{28}) a ogni fermata.
\end{enumerate}
\end{prgtry}

\key{Resume} ignora una volta il breakpoint corrente, così che ogni pressione
esegua un frame intero.

\subsection{Che cosa mostra il listato}

Il debugger non mostra il file sorgente: mostra il codice come l'assemblatore
lo ha lasciato (Listato~\ref{lst:v1:dis}).

\txtfile[caption={\code{verm1\_update}, disassemblata (ridotta; gli offset esatti dipendono dalla compilazione).},label=lst:v1:dis]{code/l07/c12.txt}

\begin{itemize}
\item I nomi sono spariti: le etichette sono diventate indirizzi, e le costanti
  sono diventate numeri.
\item Le pseudoistruzioni sono espanse: \ins{la} è \ins{auipc} più
  \ins{addi}, \ins{li} è \ins{addi} da \reg{zero}, \ins{ret} è un salto a
  \reg{ra}.
\item Alcune istruzioni cominciano con \code{c.}: sono le forme
  \emph{compresse}, a 16 bit, delle stesse istruzioni (la C di RV32IMAC).
\item La cartuccia è caricata all'indirizzo \hex{40800000}: ogni indirizzo
  mostrato parte da lì.
\end{itemize}

I formati delle istruzioni sono l'argomento del Capitolo~\ref{ch:instructions};
le forme compresse vi sono solo nominate, e sono elencate nell'Appendice~\ref{app:riscv}.

\subsection{Leggere la memoria: little-endian, dal vivo}

La vista \emph{Memory} mostra 128 byte a partire dall'indirizzo scritto, in
esadecimale (Figura~\ref{fig:v1:debug}). La variabile \code{step\_q8} è una
parola: i suoi quattro byte arrivano con il byte meno significativo per primo,
così che con $\code{step\_q8}=240$ si legge \texttt{F0 00 00 00}. Nel file
l'ordine delle variabili è \code{head\_x}, \code{head\_y}, \code{step\_q8},
quattro byte ciascuna: se \code{step\_q8} è all'indirizzo $A$, \code{head\_x} è
ad $A-8$.

Il debugger osserva soltanto. Per cambiare il programma: modificare, compilare,
eseguire.

\begin{figure}[htb]
\centering
\includegraphics[width=0.9\textwidth]{prg32qt_debugger_details.png}
\caption{\prgqt: listato, registri, memoria.}\label{fig:v1:debug}
\end{figure}

% ----------------------------------------------------------------------------
% Generated by tools/textbook/export_listings.js from tools/concepts/assessments.json: do not edit.
\begin{chapstatus}{Il filo del capitolo}
I numeri sono diventati posizioni, colori e una velocità: l'intero blocco sull'informazione, in un gioco che gira. Prossimo passo: la logica che permetterà al giocatore di guidarlo.
\par\medskip\textbf{Biglietto d'uscita} (da due a quattro minuti, da soli, su carta):
\begin{enumerate}[nosep,leftmargin=2em]
\item Verm1 dichiara title: .asciz "VERM1" e head\_x: .word 5. In quale sezione va ciascuno, e perché head\_x non deve stare in .rodata?
\item step\_q8 cresce di SPEED\_Q8 a ogni frame, e la testa si sposta di una cella ogni volta che arriva a 256. Con SPEED\_Q8 = 64, quanti frame tra due passi? Con 128? Prevedete, poi eseguite.
\end{enumerate}
\end{chapstatus}


\section{Verificate di aver capito}

\begin{quiz}{Con $\code{SPEED\_Q8}=64$, quanti frame passano tra due passi della testa?}
\item 2
\item 4
\item 64
\item 256
\end{quiz}

\begin{quiz}{Quale pixel è l'angolo in alto a sinistra della cella in colonna 10, riga 3?}
\item $(10, 3)$
\item $(80, 24)$
\item $(80, 32)$
\item $(88, 32)$
\end{quiz}

\begin{quiz}{Perché \code{verm1\_draw} salva \reg{ra} e \code{verm1\_update} no?}
\item \code{draw} è più lunga
\item \code{draw} chiama altre funzioni, che sovrascrivono \reg{ra}
\item \code{update} non usa registri
\item in \code{update} \reg{ra} è salvato automaticamente
\end{quiz}

\section{Fatelo vostro}

\begin{exercise}[su \prg]
\begin{enumerate}[nosep]
\item Fate muovere la testa verso sinistra: quale istruzione cambia, e quale
  confronto rileva il bordo?
\item Fatela muovere invece verso il basso, rientrando dalla riga 23 alla riga 0.
\item Disegnate una seconda cella, in un verde più scuro, una colonna dietro la
  testa: un primo accenno di corpo.
\item Scrivete il vostro nome accanto al titolo, in un colore che avete
  calcolato a mano in RGB565.
\end{enumerate}
\expected{Ogni modifica è di poche righe. Tenete una copia della parte~1
funzionante prima di sperimentare: la parte~2 parte da lì. (Suggerimenti: per
andare a sinistra, \code{addi t3, t3, -1}, poi \code{bgez t3} per saltare il
rientro, altrimenti \code{li t3, COLS-1}: il confronto è sul segno. Un verde
più scuro è \hex{03e0}.)}
\end{exercise}

\section{Dove siamo, e che cosa viene dopo}

Avete una cartuccia che funziona, costruita da un file scritto e capito riga
per riga. Ha uno stato in memoria, una regola che cambia lo stato a ogni frame,
e una funzione che lo disegna: la forma di ogni gioco. Manca il giocatore: la
testa va dove vuole. I Capitoli da \ref{ch:boole1} a~\ref{ch:combinational}
danno gli strumenti: bit e maschere per leggere il joystick, regole booleane,
tabelle che selezionano. La parte~2 aggiunge la guida, il cibo e un punteggio.
